\section{Preliminaries}
\label{sec:prelim}

\subsection{ARC-AGI-3 and its score}
An ARC-AGI-3 environment is a turn-based game made of a sequence of levels \citep{arcagi3}. At each step the agent
receives a frame, a grid of up to $64{\times}64$ cells with one of 16 colours each, together with the game state
(\code{NOT\_FINISHED}, \code{WIN}, \code{GAME\_OVER}) and the number of levels completed so far. It answers with one
action: \code{ACTION1}--\code{ACTION4} (directions), \code{ACTION5} (a game-specific key), \code{ACTION6}$(x,y)$ (a
click), \code{ACTION7}, or \code{RESET}. Nothing tells the agent what any action does, what it controls, or what
counts as winning. Early levels teach a mechanic and later levels compose several. The games use only core-knowledge
priors \citep{spelke2007core}: objects, agents, goals, number and elementary geometry.

The score is \emph{Relative Human Action Efficiency} (RHAE). Let $h_{l}$ be the human baseline for level $l$ (the
upper-median best human action count) and $a_{l}$ the agent's action count. Then
\begin{equation}
S_l = \min\!\Big(1.15,\, \big(\tfrac{h_l}{a_l}\big)^2\Big),\quad
E = \frac{\min\!\big(\sum_{l\le k} l,\; \sum_{l\le n} l\, S_l\big)}{\sum_{l\le n} l},
\label{eq:rhae}
\end{equation}
where $k$ of the $n$ levels were completed, and the benchmark score is the mean of $E$ over games. Three properties
drive everything that follows. Reasoning is free and only committed actions count. The penalty is \emph{quadratic},
so an agent that needs twice the human's actions keeps a quarter of the credit. Later levels weigh more, so the
mechanics learned early pay off late.

\subsection{The offline Kaggle setting}
The Kaggle track reruns a submitted notebook on hidden games, without internet access. Our measured hardware is one NVIDIA RTX PRO 6000
Blackwell (96\,GB), 176\,GB of RAM and a nine-hour limit on the GPU notebook. Every model must be open-weight and loaded
from an attached dataset. In practice the budget is fixed in GPU-hours, so what an agent can afford to think is
limited by tokens per second, and what it can afford to see is limited by the prompt length the cache keeps warm.

\subsection{How ARC has been attacked}
\paragraph{Static ARC.} The first ARC competitions were won by search over hand-written domain-specific languages
\citep{wind2020icecuber}. Later work added object graphs \citep{xu2023arga}, library learning \citep{ellis2021dreamcoder}
and minimum-description-length pressure \citep{ferre2025madil}. The ARC Prize 2024 winners moved to neural methods:
test-time training \citep{akyurek2024ttt,arcprize2024}, combined induction and transduction \citep{li2024induction},
and later tiny recursive networks \citep{jolicoeur2025trm}. These systems map a few input--output pairs to one answer.
They never have to act.

\paragraph{Agents for ARC-AGI-3.} Interactive ARC rewards a different design. General LLM agents interleave reasoning
and acting \citep{yao2022react} or grow a library of code skills \citep{wang2023voyager}. Symbolica's \emph{Arcgentica}
\citep{knutsen2026arcgentica} uses an orchestrator that never touches the game and delegates to subagents that return
compressed summaries, which keeps the context small. \emph{VISTA} \citep{vista2026} is a visual harness with lossless
frame memory. \emph{Executable World Models} \citep{rodionov2026ewm} keep a Python model of the game, verify it
against every past step, and plan through it, in the spirit of WorldCoder \citep{tang2024worldcoder} and code world
models \citep{lehrach2025cwm}. \emph{Tycho} \citep{lehmann2026tycho} induces programmatic world models. Graph-based
exploration \citep{rudakov2025graph} needs no training and ranks actions by salience. \emph{Prime Agent}
\citep{karten2026prime} gives a frontier model one persistent IPython REPL and reports a large ARC-AGI-3 gain from the
harness alone; our harness started as a port of its host loop. All of these use frontier API models. Their numbers come
from the 25 public games, which are partly memorised by now, and none of them runs offline.

\paragraph{The open-weight field.} The comparison class for this paper is the Kaggle offline track
(Table~\ref{tab:field}). Milestone~\#1 (June 2026) was won by Tufa Labs' \emph{Duck} \citep{duck2026,arcprizemilestone1}.
It runs Qwen~3.6 27B locally inside a Python REPL that exposes the game as variables, perceives through images, ASCII
grids and segmentation, and evicts its oldest messages so it can play indefinitely. Its authors found that hand-crafted
tools hurt and that letting the model improvise worked better. Second and third place both ran Gemma-4-31B as a
vision model that emits JSON actions from labelled frames \citep{arcprizemilestone1}.

\begin{table}[t]
\centering
\small
\setlength{\tabcolsep}{3.5pt}
\begin{tabular}{@{}llcr@{}}
\toprule
System & Model & Offline & Score \\
\midrule
\multicolumn{4}{@{}l}{\emph{Public games, frontier APIs (self-reported)}} \\
Exec.\ World Models & GPT-5.5 & -- & 58.1 \\
Prime Agent & frontier & -- & 95.5 \\
Arcgentica, VISTA, Tycho & frontier & -- & high \\
\midrule
\multicolumn{4}{@{}l}{\emph{Hidden games, Kaggle, one GPU}} \\
Reki (M\#1, 2nd) & Gemma-4-31B & \checkmark & 0.87 \\
Duck (M\#1, 1st) & Qwen 3.6 27B & \checkmark & 1.21 \\
\textbf{\textsc{Sarbloh} (ours)} & \textbf{Flash-Next} & \checkmark & \textbf{16.53} \\
\bottomrule
\end{tabular}
\caption{Where \textsc{Sarbloh} sits. The two halves are not comparable: one is the public set with unbounded
frontier inference, the other the hidden set on one GPU. M\#1 scores are from June. The leaderboard has moved since
(\S\ref{sec:results}).}
\label{tab:field}
\end{table}
